% SC26 MLIR tutorial -- Session 1, Part 1 (CPU). Presenter deck.
%
% Build:  make            (the deck PDF)
% See presentation/README.md.
\providecommand{\notesmode}{0}
\documentclass[aspectratio=169,10pt]{beamer}
\usepackage{sc26}
\ifnum\notesmode=1 \setbeameroption{show notes}\fi

\title{Introduction to MLIR}
\subtitle{Session 1, Part 1 --- CPU}
\author{SC26 Tutorial}
\date{}

\begin{document}

\begin{frame}[plain]
  \titlepage
\end{frame}
\note{Pre-run setup cells 0--7 of the walkthrough notebook before the room
fills: the toolchain download and the first cmake configure are the two
slowest things in the day and neither teaches anything.

Plan for the 90 minutes: 10 setup, 30 step 1, 28 defining the op,
17 writing the lowering, 5 wrap. Exercise 1 is 12 minutes of keyboard time,
exercises 2+3 another 12, exercise 4 ten. If you are behind at exercise 4,
type it yourself and narrate.}

% =======================================================================
\begin{frame}{One Computation, Two Paths}
  \begin{center}
    Today's example is \alert{matmul}:\quad
    \cd{out[i, j] += sum\_k a[i, k] * b[k, j]}
  \end{center}
  \vspace{-0.4em}
  \begin{center}
  \begin{tikzpicture}[
      node distance=0.6cm and 0.9cm,
      box/.style={draw=scblue,fill=scwash,rounded corners=2pt,inner sep=5pt,
                  font=\ttfamily\scriptsize,minimum height=2.4em,align=center},
      lbl/.style={font=\scriptsize\itshape,color=scgrey},
      arr/.style={-{Stealth[length=2mm]},draw=scgrey,thick}]

    \node[box] (scf) {scf.for nest\\\normalfont\itshape you write it};
    \node[box,right=1.6cm of scf] (llvm) {llvm dialect};
    \node[box,right=1.9cm of llvm] (run) {JIT + run};

    \draw[arr] (scf) -- node[lbl,above=2pt]{7 passes} (llvm);
    \draw[arr] (llvm) -- node[lbl,below=3pt,font=\tiny\itshape]{ExecutionEngine} (run);

    \visible<2->{
      \node[box,below=1.5cm of scf,draw=scaccent] (op)
        {blas.matmul\\\normalfont\itshape custom op you define};
    }
    \visible<3->{
      \draw[arr,draw=scaccent] (op) -- node[lbl,right,align=left]
        {your pass lowers it\\to the loop form} (scf);
      \node[lbl,below=0.35cm of run,align=center,text width=3.2cm]
        {same pipeline,\\same answer};
    }
  \end{tikzpicture}
  \end{center}

  \vspace{0.2em}
  \begin{itemize}
    \item<1-> First you write the matmul \alert{loop nest} by hand.
    \item<2-> Then you \alert{define an operation} that represents the matmul
      computation as a whole.
    \item<3-> Then you \alert{lower} from that operation to the loop form.
  \end{itemize}
\end{frame}
\note{Say: the running example today is matmul. First you write the loop nest
yourself in the scf dialect, lower it to LLVM, JIT it. Then you define an
operation that stands for the whole computation, and write the pass that
lowers it to the loop nest you just wrote.

This diagram comes back near the end, once the pass is written.}

% =======================================================================
\begin{frame}{Hands-On Exercises}
  \begin{itemize}
    \item<1-> \textbf{You edit MLIR and C++ files in Colab when prompted.}
      Edit the file, save, and re-run the cells that build and run it.
    \item<2-> \textbf{We walk through each notebook cell together}, as a
      class.
    \item<3-> \textbf{Every exercise has a complete matrix-\emph{vector}
      version next to it} --- same job, one dimension smaller. That is your
      reference.
    \item<4-> \textbf{Colab trick:} lost your place? \cd{Ctrl+F8} runs every
      cell above the one you are on.
  \end{itemize}
\end{frame}
\note{Four exercises. The matvec reference is the answer sheet for the shape,
not for the content -- say that explicitly or they will expect to copy it.
Each reference is introduced on its own exercise slide.

Ctrl+F8 is "Run before". There is no menu entry for it in the current Colab
UI, so the shortcut is the only way to say it.}

% =======================================================================
\section{Step 1 --- Mat-Mult with Explicit Loops}
\begin{frame}[plain]\sectionpage\end{frame}

\begin{frame}[fragile]{Mat-Mult with Explicit Loops}
  A matrix-matrix product that accumulates into \cd{out}:
  \begin{center}
    \cd{out[i, j] += sum\_k a[i, k] * b[k, j]}
  \end{center}
  It goes in \cd{session1/matmul\_scf.mlir}, which ships empty:
\begin{lstlisting}[language=mlir]
func.func @matmul(%a: memref<4x3xf32>, %b: memref<3x5xf32>,
                  %out: memref<4x5xf32>) {
  // your nested scf.for loops go here
  return
}
\end{lstlisting}

  \vspace{0.3em}
  \begin{itemize}
    \item<2-> Use \alert{\cd{matvec\_scf.mlir}} as your reference --- same
      directory, same computation one dimension smaller, complete.
    \item<3-> The goal: \alert{execute} \cd{@matmul} and get the expected
      answer.
  \end{itemize}

\end{frame}
\note{Switch to the ATTENDEE notebook and run its JIT check (cell 17) to show
the all-zero gate, then come back. It is worth 30 seconds: the empty function
parses and runs, so a green parse is not a passing test, and saying so stops
the first hand going up.}

% =======================================================================
\begin{frame}[fragile]{Reference: \cd{matvec\_scf.mlir}, Step 1 of 3 --- Bounds and Row Loop}
  \begin{stage}{6.0cm}
\begin{lstlisting}[language=mlir]
func.func @matvec(%mat: memref<4x3xf32>, %vec: memref<3xf32>,
                  %out: memref<4xf32>) {
  %c0 = arith.constant 0 : index
  %c1 = arith.constant 1 : index
  %numRows = memref.dim %mat, %c0 : memref<4x3xf32>
  %numCols = memref.dim %mat, %c1 : memref<4x3xf32>
  scf.for %i = %c0 to %numRows step %c1 {        // the row loop

    // the j-loop goes here: it reduces row i against %vec (next slide)

  }
  return
}
\end{lstlisting}
  \vspace{0.4em}
  \begin{itemize}
    \item<2-> \cd{memref.dim} retrieves the size of the buffer's dimensions.
    \item<3-> The sizes are fixed, 4 $\times$ 3, only to keep the example
      simple. The bounds come from \cd{memref.dim}, so the same loops work on
      \cd{memref<?x?xf32>}, whose sizes are known only at run time.
  \end{itemize}
  \end{stage}
\end{frame}

\begin{frame}[fragile]{Reference: \cd{matvec\_scf.mlir}, Step 2 of 3 --- Reduction}
  \begin{stage}{5.8cm}
\begin{lstlisting}[language=mlir]
// the j-loop: out[i] += sum_j mat[i, j] * vec[j]
(@\hl<2>{\%init = \lkw{memref.load} \%out[\%i] : \lkw{memref}<4xf32>}@)
(@\hl<4>{\%acc}@) = scf.for %j = %c0 to %numCols step %c1
                (@\hl<2>{\lkw{iter\_args}(\%a = \%init)}@) -> (f32) {
  %m = memref.load %mat[%i, %j] : memref<4x3xf32>
  %v = memref.load %vec[%j] : memref<3xf32>
  %prod = arith.mulf %m, %v : f32
  (@\hl<3>{\%a2 = \lkw{arith.addf} \%a, \%prod : \lkw{f32}}@)
  (@\hl<3>{\lkw{scf.yield} \%a2}@) : f32
}
\end{lstlisting}
  \vspace{0.4em}
  \begin{itemize}
    \item<2-> On the \alert{first iteration}, \cd{\%a} is bound to
      \cd{\%init} --- the current \cd{out[i]}, so the sum starts from what
      \cd{out} already holds.
    \item<3-> The body \alert{yields} a value, and that value is bound to
      \cd{\%a} in the \alert{next iteration}.
    \item<4-> The yield from the \alert{last iteration} is the result of the
      loop: \cd{\%acc}.
  \end{itemize}
  \end{stage}
\end{frame}

\begin{frame}[fragile]{Reference: \cd{matvec\_scf.mlir}, Step 3 of 3 --- Store}
\begin{lstlisting}[language=mlir]
func.func @matvec(%mat: memref<4x3xf32>, %vec: memref<3xf32>,
                  %out: memref<4xf32>) {
  %c0 = arith.constant 0 : index
  %c1 = arith.constant 1 : index
  %numRows = memref.dim %mat, %c0 : memref<4x3xf32>
  %numCols = memref.dim %mat, %c1 : memref<4x3xf32>
  scf.for %i = %c0 to %numRows step %c1 {                 // the row loop
    %init = memref.load %out[%i] : memref<4xf32>
    %acc = scf.for %j = %c0 to %numCols step %c1          // the j-loop
                    iter_args(%a = %init) -> (f32) {
      %m = memref.load %mat[%i, %j] : memref<4x3xf32>
      %v = memref.load %vec[%j] : memref<3xf32>
      %prod = arith.mulf %m, %v : f32
      %a2 = arith.addf %a, %prod : f32
      scf.yield %a2 : f32
    }
    (@\df{memref.store \%acc, \%out[\%i] : memref<4xf32>}@)           (@{\color{scaccent}$\leftarrow$ the value is written out}@)
  }
  return
}
\end{lstlisting}
  \vspace{0.4em}
  \begin{itemize}
    \item \cd{memref.store} writes the loop's result, \cd{\%acc}, out to
      \cd{out[i]}.
  \end{itemize}
\end{frame}
\note{Walk the file in these three steps before they start. It is 20 lines
and it contains every construct they need. The slides show the file as it is
in the repo; the // comments marking the two loops are the only additions.

Three things to land: memref.dim retrieves a dimension's size; iter\_args is
how SSA writes a reduction -- the loop carries a value, the body yields the
next one, the loop's RESULT is the final one, there is no mutable accumulator;
and the store writes that result back.}

% =======================================================================
\begin{frame}[fragile]{Exercise 1 --- Write the Loop Nest}
  \begin{center}\large
    Fill in \cd{@matmul} in \cd{session1/matmul\_scf.mlir}.
  \end{center}
  \vspace{0.2em}
  \textbf{What is different from matvec:}
  \begin{itemize}
    \item \cd{b} and \cd{out} are matrices now, so there is one more dimension
      to retrieve: the columns of \cd{b}.
    \item One more loop, over the columns of \cd{b}, between the row loop and
      the reduction loop.
    \item \cd{out} and \cd{b} take two indices: \cd{out[i, j]} and
      \cd{b[k, j]}.
  \end{itemize}
  \vspace{0.5em}
  \textbf{Check as you go} --- after every save, re-run the round-trip cell:
\begin{lstlisting}[style=plain,basicstyle=\ttfamily\small]
!mlir-opt matmul_scf.mlir
\end{lstlisting}
  \vspace{0.2em}
  \begin{itemize}
    \item No pass pipeline: \cd{mlir-opt} parses and re-prints, nothing else.
    \item A clean round-trip means it \alert{parses and type-checks}, not
      that it is right. We verify the answer together in the solution.
  \end{itemize}
  \vspace{0.5em}
  \begin{center}\color{scaccent}\large
    12 minutes to write it --- then we walk through the solution.
  \end{center}
\end{frame}
\note{Walk around while they type -- the two mistakes you will see most are
starting the reduction from a zero constant, and indexing b as [j, k]. Neither
produces an error, which is why the round-trip is only the immediate check.

Worth saying about the round-trip: it is the most useful habit in MLIR work.
Before you debug a pipeline, round-trip the input. The printed form may differ
from what they typed -- generated SSA names, normalised whitespace. That is
the printer, not a transformation.}

% =======================================================================
\begin{frame}[fragile]{Solution --- \cd{matmul\_scf.mlir}}
\begin{lstlisting}[language=mlir]
func.func @matmul(%a: memref<4x3xf32>, %b: memref<3x5xf32>,
                  %out: memref<4x5xf32>) {
  %c0 = arith.constant 0 : index
  %c1 = arith.constant 1 : index
  (@\hl<2>{\%M = \lkw{memref.dim} \%a, \%c0 : \lkw{memref}<4x3xf32>}@)
  (@\hl<2>{\%K = \lkw{memref.dim} \%a, \%c1 : \lkw{memref}<4x3xf32>}@)                 (@\only<2>{\color{scaccent}$\leftarrow$ find M, N and K}@)
  (@\hl<2>{\%N = \lkw{memref.dim} \%b, \%c1 : \lkw{memref}<3x5xf32>}@)
  (@\hl<3>{\lkw{scf.for} \%i = \%c0 \lkw{to} \%M \lkw{step} \%c1 \{}@)                         (@\only<3>{\color{scaccent}$\leftarrow$ loop over rows and columns}@)
    (@\hl<3>{\lkw{scf.for} \%j = \%c0 \lkw{to} \%N \lkw{step} \%c1 \{}@)
      (@\hl<4>{\%init = \lkw{memref.load} \%out[\%i, \%j] : \lkw{memref}<4x5xf32>}@)
      %acc = scf.for %k = %c0 to %K step %c1
                      iter_args(%s = %init) -> (f32) {
        (@\hl<4>{\%av = \lkw{memref.load} \%a[\%i, \%k] : \lkw{memref}<4x3xf32>}@)      (@\only<4>{\color{scaccent}$\leftarrow$ read a and b}@)
        (@\hl<4>{\%bv = \lkw{memref.load} \%b[\%k, \%j] : \lkw{memref}<3x5xf32>}@)
        %prod = arith.mulf %av, %bv : f32
        %s2 = arith.addf %s, %prod : f32
        scf.yield %s2 : f32
      }
      (@\hl<4>{\lkw{memref.store} \%acc, \%out[\%i, \%j] : \lkw{memref}<4x5xf32>}@) (@\only<4>{\color{scaccent}$\leftarrow$ write the accumulated value to out}@)
    }
  }
  return
}
\end{lstlisting}
\end{frame}
\note{This is where they verify their own answer. Run the pipeline and the JIT
check in the walkthrough notebook as you go.

Three things to say that are not on the slide. At the loads: it is
\%b[\%k, \%j], the reduction index first -- swap them and nothing complains,
you just get the wrong answer. At the dims: \%M and \%N are the two parallel
dimensions and \%K is the reduction one; that distinction is the whole of
Session 2. And at the end: remember this shape, they build it again in C++.}

% =======================================================================
\begin{frame}[fragile]{The Lowering Pipeline}
\begin{lstlisting}[style=plain,basicstyle=\ttfamily\scriptsize]
func.func(llvm-request-c-wrappers)   (@\only<2>{\color{scaccent}$\leftarrow$ a C-callable @matmul wrapper, for Python}@)
convert-scf-to-cf                    (@\only<3>{\color{scaccent}$\leftarrow$ structured loops become a CFG}@)
convert-cf-to-llvm                   (@\only<4>{\color{scaccent}$\leftarrow$ branches become llvm.br / llvm.cond\_br}@)
convert-arith-to-llvm                (@\only<5>{\color{scaccent}$\leftarrow$ mulf, addf}@)
finalize-memref-to-llvm              (@\only<6>{\color{scaccent}$\leftarrow$ loads become getelementptr + load}@)
convert-func-to-llvm                 (@\only<7>{\color{scaccent}$\leftarrow$ signatures, calling convention}@)
reconcile-unrealized-casts           (@\only<8>{\color{scaccent}$\leftarrow$ clean up the type-conversion seams}@)
\end{lstlisting}
  \vspace{0.25em}
  \begin{itemize}
    \item<1-> Seven passes, each lowering \alert{one} dialect.
    \item<9-> After every one of them you have a valid MLIR module you can
      print and read.
  \end{itemize}
  \visible<9->{\cmdbox{mlir-opt ... -{}-mlir-print-ir-after-all -{}-mlir-disable-threading}}
  \begin{itemize}
    \item<10-> The result is MLIR's \cd{llvm} \emph{dialect} --- not LLVM IR.
      Translation to real LLVM IR happens inside the execution engine.
  \end{itemize}
\end{frame}
\note{This is the conceptual centre of Part 1. Do not rush it.

If someone asks why it is this many passes: because each one is independently
testable and reusable, and the GPU path in Part 2 reuses most of them.

Scroll the real lowered output in the notebook: point at one llvm.cond\_br
(the loops became a CFG) and one getelementptr chain (the loads became
explicit address arithmetic), then move on.

The two flags are mlir-opt flags. The notebook cell runs this pipeline
through Python's PassManager, so to show them, run mlir-opt on the file with
the same pipeline string:
  mlir-opt matmul\_scf.mlir -{}-pass-pipeline="builtin.module(...)"
      -{}-mlir-print-ir-after-all -{}-mlir-disable-threading
Threading: with it on, function passes run on several functions at once and
their dumps interleave.}

% =======================================================================
\begin{frame}[fragile]{Execution}
\begin{lstlisting}[language=Python,basicstyle=\ttfamily\scriptsize,keywordstyle=\color{scblue}]
a = np.arange(12, dtype=np.float32).reshape(4, 3)
b = np.arange(15, dtype=np.float32).reshape(3, 5)
out = np.zeros((4, 5), dtype=np.float32)

run_and_check(module, "matmul", [a, b, out], expected=a @ b)
\end{lstlisting}
  \vspace{0.5em}
  \begin{itemize}
    \item<2-> numpy arrays are passed as \alert{memref descriptors} --- a
      struct of pointer, offset, sizes and strides. No copying.
    \item<3-> The last argument is the output buffer, by convention.
      \cd{run\_and\_check} \alert{zeroes it again} before every run, because
      the notebook reuses it from step to step.
    \item<4-> Checked against numpy's own matmul.
  \end{itemize}
  \vspace{0.4em}
  \visible<5->{\begin{center}\color{scgreen}\large\ttfamily OK\end{center}}
\end{frame}
\note{If someone's run gives all zeros while mlir-opt was happy, they saved
the file but did not re-run from the round-trip cell.}

% =======================================================================
\begin{frame}{Step 1 Recap}
  \begin{center}
    \large You wrote a matrix multiply at the loop level,\\
    compiled it through seven passes, and ran it.
  \end{center}
  \vspace{1em}
  \begin{itemize}
    \item<2-> \cd{scf.for} with \cd{iter\_args} for the reduction.
    \item<3-> \cd{memref.dim} so the bounds come from the buffer.
    \item<4-> One pipeline, dialect by dialect, down to \cd{llvm}.
  \end{itemize}
  \vspace{1em}
  \visible<5->{\begin{center}\color{scaccent}\large
    Short break. If anything is not working for you, raise your hand.
  \end{center}}
\end{frame}

% =======================================================================
\section{Step 2 --- BLAS Dialect}
\begin{frame}[plain]\sectionpage\end{frame}

\begin{frame}[fragile]{Defining Your Own Level of Abstraction}
  \begin{columns}[t]
    \begin{column}{0.47\textwidth}
      \begin{center}\small\textbf{the loop nest}\end{center}
\begin{lstlisting}[language=mlir,basicstyle=\ttfamily\scriptsize]
scf.for %i = ... {
  scf.for %j = ... {
    %0 = memref.load %out[%i, %j]
    %1 = scf.for %k = ... iter_args(...) {
      ...
    }
    memref.store %1, %out[%i, %j]
  }
}
\end{lstlisting}
      \begin{center}\scriptsize\color{scgrey}
        20 operations.
      \end{center}
    \end{column}
    \begin{column}{0.47\textwidth}
      \onslide<2->
      \begin{center}\small\textbf{an operation that represents it}\end{center}
\begin{lstlisting}[language=mlir,basicstyle=\ttfamily\scriptsize]
blas.matmul %a, %b, %out
  : memref<4x3xf32>,
    memref<3x5xf32>,
    memref<4x5xf32>
\end{lstlisting}
      \begin{center}\scriptsize\color{scgrey}
        One operation.\end{center}
      \onslide<1->
    \end{column}
  \end{columns}
  \vspace{0.5em}
  \begin{itemize}
    \item<3-> MLIR lets you create \alert{higher levels of abstraction} ---
      your own DSL --- as operations of your own.
    \item<4-> The compiler then reasons about \alert{a matmul operation}
      directly, instead of deriving that from a loop nest by pattern
      matching. Nothing has to be lifted back out of the loops.
    \item<5-> One definition, \alert{many lowerings}: CPU loops in this part,
      \cd{gpu.launch} in Part 2.
  \end{itemize}
\end{frame}
\note{This is the argument for the whole session, so do not skip it. The op
is not shorter -- it is more \emph{informative}. A tile, a distribute, a fuse
becomes a rewrite on one op rather than a loop-dependence analysis, and
everything Session 2 does with the transform dialect is only possible because
the op says what it is.}

% =======================================================================
\begin{frame}{Anatomy of a Dialect}
  \begin{center}
  \begin{tikzpicture}[
      box/.style={draw=scblue,fill=scwash,rounded corners=2pt,inner sep=6pt,
                  font=\ttfamily\scriptsize,align=center,minimum width=3.2cm},
      lbl/.style={font=\scriptsize,color=scink,align=center},
      arr/.style={-{Stealth[length=2mm]},draw=scgrey,thick}]

    \node[box] (ods) {BlasOps.td.txt};
    \node[lbl,below=0.15cm of ods,text width=3.8cm]
      {\textbf{declarative}\\operands, types, syntax};

    \node[box,right=1.1cm of ods] (ver) {BlasOps.cpp};
    \node[lbl,below=0.15cm of ver,text width=3.8cm]
      {\textbf{imperative}\\what ODS cannot say};

    \draw[arr] (ods) -- (ver);

    \visible<2->{
      \node[box,right=1.1cm of ver,draw=scaccent] (low) {LowerBlasToSCF.cpp};
      \node[lbl,below=0.15cm of low,text width=3.8cm]
        {\textbf{meaning}\\one pattern per op};
      \draw[arr] (ver) -- (low);}

    \visible<3->{
      \draw[decorate,decoration={brace,amplitude=5pt,mirror},draw=scgreen]
        ($(ods.south west)+(0,-1.25)$) -- ($(low.south east)+(0,-1.25)$);
      \node[draw=scgreen,fill=white,rounded corners=2pt,inner sep=5pt,
            font=\ttfamily\scriptsize] (so)
        at ($(ver.south)+(0,-2.3)$) {libBlasPlugin.so};
      \node[lbl,right=0.4cm of so,text width=4.6cm,align=left]
        {one shared object, loaded by \cd{mlir-opt} at run time:\\
         \cd{-{}-load-dialect-plugin}, \cd{-{}-load-pass-plugin}};}
  \end{tikzpicture}
  \end{center}
  \vspace{0.5em}
  \begin{itemize}
    \item<1-> First you \alert{define the op}. Its \alert{ODS} (Operation
      Definition Specification) is a declarative spec written in TableGen.
      \cd{mlir-tblgen} generates the op's C++ from it, and the verifier in
      \cd{BlasOps.cpp} adds what ODS cannot say.
    \item<2-> Then you \alert{write the lowering} that turns it into loops.
  \end{itemize}
\end{frame}
\note{TableGen files are normally .td. Colab's editor will not open that
extension, so the one file attendees edit is BlasOps.td.txt; mlir-tblgen does
not care about the extension. The dialect's other .td files keep their real
names, since nobody needs to open them.}

% =======================================================================
\begin{frame}[fragile]{Defining an Operation}
  We want to define an operation that looks like this:
\begin{lstlisting}[language=mlir]
func.func @matmul(%a: memref<4x3xf32>, %b: memref<3x5xf32>,
                  %out: memref<4x5xf32>) {
  blas.matmul %a, %b, %out
    : memref<4x3xf32>, memref<3x5xf32>, memref<4x5xf32>
  return
}
\end{lstlisting}
  \vspace{0.25em}
  \begin{itemize}
    \item<2-> That is \cd{session1/matmul\_custom.mlir} --- the MLIR format we
      expect is in that file, already written.
    \item<3-> To start with, round-tripping it \alert{fails}: the operation is
      not defined yet.
  \end{itemize}
  \visible<3->{\errbox{matmul\_custom.mlir:5:3: error: custom op 'blas.matmul' is unknown}}
  \vspace{0.3em}
  \begin{itemize}
    \item<4-> Once you have defined the operation, the same round-trip
      \alert{succeeds}: \cd{mlir-opt} parses the file and prints it back
      unchanged.
  \end{itemize}
\end{frame}
\note{Parsing the file proves their assemblyFormat accepts the syntax;
printing it back identically proves it produces the same syntax. One check,
both directions -- which is what the assembly-format step is about.

The real file has three comment lines above the function, which is why the
error says line 5.}

% =======================================================================
\begin{frame}[fragile]{Reference: \cd{Blas\_MatVecOp}, Step 1 of 5 --- Operands}
  \cd{blas-dialect/include/Blas/BlasOps.td.txt}:
\begin{lstlisting}[language=ods]
def Blas_MatVecOp : Blas_Op<"matvec"> {
  let summary = "matrix-vector product: out += mat * vec";

  let arguments = (ins
    MemRefRankOf<[F32], [2]>:$mat,           (@\only<3>{\color{scaccent}$\leftarrow$ a 2D f32 memref, named \cd{mat}}@)
    MemRefRankOf<[F32], [1]>:$vec,
    MemRefRankOf<[F32], [1]>:$out);

  // steps 2 and 4: assemblyFormat and hasVerifier
}
\end{lstlisting}
  \vspace{0.4em}
  \begin{itemize}
    \item<2-> \cd{Blas\_Op<"matvec">} makes this the op \cd{blas.matvec}. Its
      C++ class is the def name minus the \cd{Blas\_} prefix:
      \cd{blas::MatVecOp}.
    \item<3-> \cd{arguments} lists the operands. Each one has a
      \alert{type constraint} and a \alert{name}.
    \item<4-> ODS generates the check for those types: a \cd{blas.matvec}
      with a 1D \cd{mat} is rejected without you writing any C++.
  \end{itemize}
\end{frame}

\begin{frame}[fragile]{Reference: \cd{Blas\_MatVecOp}, Step 2 of 5 --- Assembly Format}
\begin{lstlisting}[language=ods]
let assemblyFormat =
  "$mat `,` $vec `,` $out attr-dict `:` type($mat) `,` type($vec) `,` type($out)";
\end{lstlisting}
  \vspace{0.3em}
  \begin{center}
  \begin{tikzpicture}[
      font=\ttfamily\scriptsize,
      piece/.style={inner sep=2pt},
      arr/.style={-{Stealth[length=1.4mm]},draw=scaccent}]

    % what mlir-opt prints
    \node[piece] (pa) at (0,0) {\%mat};
    \node[piece,right=0.1cm of pa] (pc1) {,};
    \node[piece,right=0.1cm of pc1] (pb) {\%vec};
    \node[piece,right=0.1cm of pb] (pc2) {,};
    \node[piece,right=0.1cm of pc2] (pout) {\%out};
    \node[piece,right=0.3cm of pout] (pcolon) {:};
    \node[piece,right=0.3cm of pcolon] (pta) {memref<4x3xf32>};
    \node[piece,right=0.1cm of pta] (pc3) {,};
    \node[piece,right=0.1cm of pc3] (ptb) {memref<3xf32>};
    \node[piece,right=0.1cm of ptb] (pc4) {,};
    \node[piece,right=0.1cm of pc4] (ptc) {memref<4xf32>};
    \node[piece,left=0.3cm of pa] (pname) {blas.matvec};

    % the format string, each piece above what it produces
    \node[piece,above=1.1cm of pa] (fa) {\color{scblue}\$mat};
    \node[piece,above=1.1cm of pb] (fb) {\color{scblue}\$vec};
    \node[piece,above=1.1cm of pout] (fout) {\color{scblue}\$out};
    \node[piece,above=1.1cm of pcolon] (fcolon) {`:`};
    \node[piece,above=1.1cm of pta] (fta) {\color{scaccent}type(\$mat)};
    \node[piece,above=1.1cm of ptb] (ftb) {\color{scaccent}type(\$vec)};
    \node[piece,above=1.1cm of ptc] (ftc) {\color{scaccent}type(\$out)};

    \visible<2->{\draw[arr] (fa) -- (pa);}
    \visible<3->{\draw[arr] (fb) -- (pb);
                  \draw[arr] (fout) -- (pout);}
    \visible<4->{\draw[arr] (fta) -- (pta);
                  \draw[arr] (ftb) -- (ptb);
                  \draw[arr] (ftc) -- (ptc);}
  \end{tikzpicture}
  \end{center}
  \vspace{0.3em}
  \begin{itemize}
    \item<5-> MLIR, like LLVM, calls the textual form of IR its
      \alert{assembly}: this string is the op's syntax in a \cd{.mlir} file.
    \item<6-> From it, TableGen generates both the \alert{parser} (text to op)
      and the \alert{printer} (op to text), so the two always agree.
    \item<7-> \cd{attr-dict} before the colon is \alert{not optional
      boilerplate}. Leave it out and your op cannot carry attributes, which
      breaks things later in ways that look unrelated.
  \end{itemize}
\end{frame}
\note{Spend a moment on attr-dict. Every year someone leaves it out and the
failure lands three steps away.}

\begin{frame}[fragile]{Reference: \cd{Blas\_MatVecOp}, Step 3 of 5 --- What TableGen Generates}
  From the ODS, \cd{mlir-tblgen} writes \cd{build/include/Blas/BlasOps.h.inc}:
\begin{lstlisting}[language=C++,basicstyle=\ttfamily\scriptsize,morekeywords={class,static,void,const}]
class MatVecOp : public ::mlir::Op<MatVecOp, ...> {
  ::mlir::TypedValue<::mlir::MemRefType> getMat();       (@\only<2>{\color{scaccent}$\leftarrow$ from \$mat}@)
  ::mlir::TypedValue<::mlir::MemRefType> getVec();       (@\only<2>{\color{scaccent}$\leftarrow$ from \$vec}@)
  ::mlir::TypedValue<::mlir::MemRefType> getOut();       (@\only<2>{\color{scaccent}$\leftarrow$ from \$out}@)
  static void build(::mlir::OpBuilder &odsBuilder, ::mlir::OperationState &odsState,
                    ::mlir::Value mat, ::mlir::Value vec, ::mlir::Value out);
                               (@\only<3>{\color{scaccent}$\leftarrow$ what a pattern calls to CREATE the op}@)
  ::llvm::LogicalResult verify();
                               (@\only<4>{\color{scaccent}$\leftarrow$ declared here, defined by you (step 5)}@)
};
\end{lstlisting}
  \vspace{0.25em}
  \begin{itemize}
    \item<2-> Every operand you name in \cd{arguments} gets a \cd{get*}
      accessor: \cd{\$mat} becomes \cd{getMat()}, \cd{\$out} becomes
      \cd{getOut()}.
    \item<5-> ODS generates the boilerplate --- accessors, builders, parser,
      printer --- and keeps operations, dialects and types \alert{consistent}
      across MLIR.
    \item<6-> Nothing in this file is edited by hand. It is regenerated on
      every build.
  \end{itemize}
\end{frame}
\note{The accessors are the point of this slide: the verifier in step 5 and
the lowering pattern later both call getMat(), getVec() and getOut(), and
this is where they come from. Their own op will get getA(), getB(), getOut().

These are the real declarations, with the long trait list on the class
elided. The notebook prints the same lines for MatMulOp after the build.}

\begin{frame}{Reference: \cd{Blas\_MatVecOp}, Step 4 of 5 --- What ODS Cannot Say}
  \begin{itemize}
    \item<1-> ODS can say ``\alert{a 2D f32 memref and two 1D ones}''.
    \item<2-> ODS \alert{cannot} say ``\cd{mat}'s column count equals
      \cd{vec}'s size''. Shape \emph{relationships} go in the C++
      \alert{\cd{verify()}} method:
  \end{itemize}
  \vspace{0.3em}
  \visible<2->{%
  \begin{center}
  \begin{tabular}{@{}ll@{}}
    \small\textbf{constraint} & \small\textbf{why}\\[0.2em]
    \cd{mat.dim(1) == vec.dim(0)} & the reduction dimension\\
    \cd{mat.dim(0) == out.dim(0)} & rows\\
  \end{tabular}
  \end{center}}
  \vspace{0.3em}
  \begin{itemize}
    \item<3-> \cd{let hasVerifier = 1;} in the ODS is what generates the
      \cd{verify()} \emph{declaration}. Leave out the C++ definition and
      \cd{mlir-opt} fails to load the plugin with an undefined-symbol error.
    \item<4-> The ODS checks plus your \cd{verify()} are the op's
      \alert{verifier}. MLIR runs it when IR is \alert{parsed} and
      \alert{after every pass}, so the op's guarantees hold on the input and
      across every transformation.
  \end{itemize}
\end{frame}
\note{If someone forgets the definition for their own op, the plugin still
builds -- we link with -Wl,-z,undefs because MLIR's own symbols are resolved
by mlir-opt at dlopen time. The failure comes when mlir-opt loads it:
  symbol lookup error: libBlasPlugin.so: undefined symbol:
  \_ZN4blas8MatMulOp6verifyEv
which demangles (c++filt) to blas::MatMulOp::verify().}

\begin{frame}[fragile]{Reference: \cd{Blas\_MatVecOp}, Step 5 of 5 --- \cd{MatVecOp::verify()}}
  \cd{blas-dialect/lib/BlasOps.cpp}:
\begin{lstlisting}[language=C++,morekeywords={LogicalResult,auto,return},basicstyle=\ttfamily\scriptsize]
LogicalResult MatVecOp::verify() {
  auto matType = cast<MemRefType>(getMat().getType());         (@\only<2>{\color{scaccent}$\leftarrow$ the generated accessors}@)
  auto vecType = cast<MemRefType>(getVec().getType());
  auto outType = cast<MemRefType>(getOut().getType());

  if (failed(verifyDimsMatch(*this, matType.getDimSize(1),     (@\only<3>{\color{scaccent}$\leftarrow$ the reduction dimension}@)
                             vecType.getDimSize(0),
                             "matrix column count must match input vector size")))
    return failure();

  return verifyDimsMatch(*this, matType.getDimSize(0), outType.getDimSize(0),   (@\only<4>{\color{scaccent}$\leftarrow$ rows}@)
                         "matrix row count must match output vector size");
}
\end{lstlisting}
  \vspace{0.3em}
  \begin{itemize}
    \item<2-> \cd{getMat()}, \cd{getVec()} and \cd{getOut()} are the
      accessors from step 3.
    \item<3-> One \cd{verifyDimsMatch} call per constraint. The helper lets a
      dynamic (\cd{?}) dimension through: its size is known only at run time,
      so the verifier cannot check it.
    \item<5-> The string is \alert{your} error message --- what someone sees
      when they write a \cd{blas.matvec} whose shapes do not agree.
  \end{itemize}
\end{frame}

% =======================================================================
\begin{frame}{Exercises 2 and 3 --- Define the Op}
  \textbf{2. The ODS}\quad \cd{blas-dialect/include/Blas/BlasOps.td.txt}
  \begin{itemize}
    \item Define \cd{Blas\_MatMulOp} where the \cd{TODO} is, modelled on
      \cd{Blas\_MatVecOp} right above it.
    \item Three 2D \cd{f32} memrefs, named \cd{\$a}, \cd{\$b} and \cd{\$out}.
      Include \mbox{\cd{let hasVerifier = 1;}}
  \end{itemize}
  \vspace{0.8em}
  \textbf{3. The verifier}\quad \cd{blas-dialect/lib/BlasOps.cpp}
  \begin{itemize}
    \item Write \mbox{\cd{LogicalResult MatMulOp::verify()}}, modelled on
      \mbox{\cd{MatVecOp::verify()}}.
    \item Three constraints, one \cd{verifyDimsMatch} call each.
  \end{itemize}
  \vspace{0.8em}
  \begin{center}\small\color{scgrey}
    These two go together: \cd{hasVerifier = 1} declares \cd{verify()},
    and exercise 3 defines it.
  \end{center}
  \vspace{0.4em}
  \begin{center}\color{scaccent}\large
    12 minutes to write them --- then we walk through the solutions.
  \end{center}
\end{frame}

% =======================================================================
\begin{frame}[fragile]{Solution --- The ODS, and What Changed from Matvec}
  \begin{stage}{5.8cm}
\begin{lstlisting}[language=ods,deletestring={[b]"}]
def Blas_(@\df{MatMulOp}@) : Blas_Op<"(@\df{matmul}@)"> {                  (@\only<2>{\color{scaccent}$\leftarrow$ a new name}@)
  let summary = "matrix-matrix product: out += a * b";

  let arguments = (ins
    MemRefRankOf<[F32], [2]>:(@\df{\$a}@),                   (@\only<3>{\color{scaccent}$\leftarrow$ new operand names}@)
    MemRefRankOf<[F32], [(@\df{2}@)]>:(@\df{\$b}@),                   (@\only<4>{\color{scaccent}$\leftarrow$ \cd{b} is a matrix now ...}@)
    MemRefRankOf<[F32], [(@\df{2}@)]>:$out);                (@\only<4>{\color{scaccent}$\leftarrow$ ... and so is \cd{out}}@)

  let assemblyFormat =                           (@\only<3>{\color{scaccent}$\leftarrow$ the same format, with the new names}@)
    "(@\df{\$a}@) `,` (@\df{\$b}@) `,` $out attr-dict `:` type((@\df{\$a}@)) `,` type((@\df{\$b}@)) `,` type($out)";

  let hasVerifier = 1;
}
\end{lstlisting}
  \vspace{0.4em}
  \begin{itemize}
    \item<5-> Copy \cd{Blas\_MatVecOp}, rename it and its operands, and make
      \cd{b} and \cd{out} 2D. That is all it is.
    \item<6-> The generated C++ class is \cd{blas::MatMulOp}, with
      \cd{getA()}, \cd{getB()} and \cd{getOut()} --- the names the lowering
      expects.
  \end{itemize}
  \end{stage}
\end{frame}

\begin{frame}[fragile]{Solution --- The Verifier, and What Changed from Matvec}
\begin{lstlisting}[language=C++,morekeywords={LogicalResult,auto,return},basicstyle=\ttfamily\scriptsize,stringstyle=\color{scink}]
LogicalResult (@\df{MatMulOp}@)::verify() {
  auto aType = cast<MemRefType>(getA().getType());
  auto bType = cast<MemRefType>(getB().getType());
  auto outType = cast<MemRefType>(getOut().getType());

  if (failed(verifyDimsMatch(*this, aType.getDimSize(1), bType.getDimSize(0),
                             "inner dimensions must match")))
    return failure();
  if (failed(verifyDimsMatch(*this, aType.getDimSize(0), outType.getDimSize(0),
                             "row count of 'a' must match row count of 'out'")))
    return failure();
  (@\df{return verifyDimsMatch(*this, bType.getDimSize(1), outType.getDimSize(1),}@)
                         (@\df{"column count of 'b' must match column count of 'out'");}@)
}
\end{lstlisting}
  \vspace{0.3em}
  \begin{center}
  \begin{tabular}{@{}lll@{}}
    \small\textbf{constraint} & \small\textbf{why} & \\[0.2em]
    \cd{a.dim(1) == b.dim(0)} & the reduction dimension & \small\color{scgrey}as in matvec\\
    \cd{a.dim(0) == out.dim(0)} & rows & \small\color{scgrey}as in matvec\\
    \textcolor{scaccent}{\cd{b.dim(1) == out.dim(1)}} & \textcolor{scaccent}{columns} & \small\textcolor{scaccent}{new}\\
  \end{tabular}
  \end{center}
\end{frame}
\note{Same shape as MatVecOp::verify() with the accessors renamed, plus one
more constraint: the output has columns now, and they come from b.

Optional demo, and it lands better than any explanation: after the
round-trip check passes, edit matmul\_custom.mlir so b is memref<4x5xf32>,
re-run, and show the room the error message THEY wrote coming out of the
compiler. Undo it afterwards.}

% =======================================================================
\begin{frame}{What the Build Does}
  \begin{center}
  \begin{tikzpicture}[
      box/.style={draw=scblue,fill=scwash,rounded corners=2pt,inner sep=4pt,
                  font=\ttfamily\scriptsize,align=center},
      gen/.style={box,draw=scgrey,fill=white},
      arr/.style={-{Stealth[length=2mm]},draw=scgrey},
      lbl/.style={font=\scriptsize\itshape,color=scgrey}]

    \node[box] (td) at (0,1.1)
      {BlasOps.td.txt\\[-0.1em]\normalfont\itshape\scriptsize you wrote this};
    \visible<3->{
      \node[box] (cpp) at (0,-1.1)
        {BlasOps.cpp, LowerBlasToSCF.cpp\\[-0.1em]%
         \normalfont\itshape\scriptsize you wrote these};}
    \visible<2->{
      \node[gen] (inc) at (5.0,1.1) {BlasOps.h.inc\\[-0.1em]BlasOps.cpp.inc};
      \draw[arr] (td) -- node[lbl,above]{mlir-tblgen} (inc);}
    \visible<4->{
      \node[gen,draw=scgreen] (so) at (9.4,0) {libBlasPlugin.so};
      \draw[arr] (inc) -- (so);
      \draw[arr] (cpp) -- node[lbl,below,pos=0.55]{clang++} (so);}
    \visible<5->{
      \node[box,draw=scaccent] (opt) at (9.4,-1.8) {mlir-opt};
      \draw[arr,draw=scaccent] (so) -- node[lbl,right]{dlopen} (opt);}
  \end{tikzpicture}
  \end{center}
  \vspace{0.5em}
  \begin{itemize}
    \item<6-> \cd{.td.txt} and \cd{.cpp} are \alert{compiled into} the plugin.
      \cd{.mlir} files are read \alert{fresh every time}.
    \item<7-> So: edit MLIR, re-run the check. Edit C++ or TableGen,
      \alert{re-run the build cell first}.
    \item<8-> When a change seems to have no effect, this is why. Every time.
  \end{itemize}
\end{frame}
\note{Say the last two bullets out loud before the first build, not after.
This is the single most common lost-ten-minutes in the session.

While it builds (1-2 min first time): ninja is running mlir-tblgen over their
.td.txt to generate a class declaration, a parser, a printer and builders, then
compiling that together with the hand-written .cpp into one shared object.

If the build fails: read the FIRST error, not the last. A TableGen syntax error
cascades into dozens of C++ errors about a class that was never generated.}

% =======================================================================
\begin{frame}[fragile]{Check 2 --- Round-Trip through the Plugin}
\begin{lstlisting}[style=plain,basicstyle=\ttfamily\small]
!mlir-opt --load-dialect-plugin=libBlasPlugin.so matmul_custom.mlir
\end{lstlisting}
  \vspace{0.5em}
  \begin{itemize}
    \item<2-> Python's \cd{Context} knows nothing about our plugin, so
      anything involving \cd{blas} shells out to \cd{mlir-opt}.
    \item<3-> Success looks like \alert{the file, printed back unchanged}.
  \end{itemize}
  \vspace{0.5em}
  \visible<4->{\errbox{matmul\_custom.mlir:5:3: error: custom op 'blas.matmul' is unknown}}
  \visible<4->{\vspace{0.3em}
  \begin{center}\small
    Either the op is not defined --- or \alert{the plugin did not rebuild}.
    In a room this size, it is the rebuild.
  \end{center}}
\end{frame}

% =======================================================================
\begin{frame}[fragile]{Reference: \cd{MatVecToSCFPattern}, Step 1 of 5 --- Scaffolding}
  \begin{columns}[t]
    \begin{column}{0.59\textwidth}
      \begin{center}\small\textbf{the C++}\end{center}
\begin{lstlisting}[style=halfs,language=C++,morekeywords={LogicalResult,Value,Location,override,struct,public,using}]
(@\hl<2>{\lkw{struct} MatVecToSCFPattern : \lkw{public} OpRewritePattern<MatVecOp> \{}@)
  using OpRewritePattern::OpRewritePattern;
  LogicalResult
  matchAndRewrite(MatVecOp op,
                  PatternRewriter &rewriter) const override {
    Location loc = op.getLoc();
    Value mat = op.getMat();

    (@\hl<3>{\lkw{Value} c0 = arith::ConstantIndexOp::create(rewriter, loc, 0);}@)
    (@\hl<3>{\lkw{Value} c1 = arith::ConstantIndexOp::create(rewriter, loc, 1);}@)
    (@\hl<4>{\lkw{Value} numRows = memref::DimOp::create(rewriter, loc, mat, c0);}@)
    (@\hl<4>{\lkw{Value} numCols = memref::DimOp::create(rewriter, loc, mat, c1);}@)
    // steps 2-4 go here
    rewriter.eraseOp(op);
    return success();
  }
};
\end{lstlisting}
    \end{column}
    \begin{column}{0.40\textwidth}
      \begin{center}\small\textbf{the IR it generates}\end{center}
\begin{lstlisting}[style=halfs,language=mlir]
(@\strut@)







(@\visible<3->{\hl<3>{\%c0 = \lkw{arith.constant} 0 : \lkw{index}}}@)
(@\visible<3->{\hl<3>{\%c1 = \lkw{arith.constant} 1 : \lkw{index}}}@)
(@\visible<4->{\hl<4>{\%numRows = \lkw{memref.dim} \%mat, \%c0}}@)
(@\visible<4->{\hl<4>{\%numCols = \lkw{memref.dim} \%mat, \%c1}}@)
\end{lstlisting}
    \end{column}
  \end{columns}
  \vspace{0.3em}
  \begin{itemize}
    \item<2-> A pattern \alert{matches one op}:
      \cd{OpRewritePattern<MatVecOp>} rewrites \cd{blas.matvec} and nothing
      else.
    \item<3-> First the index constants, \cd{0} and \cd{1}, for the loop
      bounds and steps.
    \item<4-> Then the dimensions of the matrix, with \cd{memref.dim}:
      \cd{c0} and \cd{c1} say which dimension.
  \end{itemize}
\end{frame}
\note{\cd{blas-dialect/lib/LowerBlasToSCF.cpp}. This pattern is complete in
the file and is what they copy from. The operands vec and out are read here
too in the file; they are left off this slide and come into use in steps 3
and 4.

On the right-hand columns of these slides: the IR each C++ call builds, with
the SSA names tidied and the types left off. It is the hand-written
matvec\_scf.mlir from step 1, rebuilt in C++.

If someone runs the pass and compares: the printed output shows
arith.constant 4 and 3 where this slide shows memref.dim. The rewrite driver
folds memref.dim of a static shape to a constant.}

\begin{frame}[fragile]{Reference: \cd{MatVecToSCFPattern}, Step 2 of 5 --- Row Loop}
  \begin{columns}[t]
    \begin{column}{0.59\textwidth}
      \begin{center}\small\textbf{the C++}\end{center}
\begin{lstlisting}[style=halfs,language=C++,morekeywords={auto,Value}]
(@\hl<2>{\lkw{auto} rowLoop = scf::ForOp::create(rewriter, loc, c0, numRows, c1);}@)
{
  (@\hl<4>{OpBuilder::InsertionGuard guard(rewriter);}@)
  (@\hl<3>{rewriter.setInsertionPointToStart(rowLoop.getBody());}@)
  Value i = rowLoop.getInductionVar();

  // steps 3 and 4 go here
}
\end{lstlisting}
    \end{column}
    \begin{column}{0.40\textwidth}
      \begin{center}\small\textbf{the IR it generates}\end{center}
\begin{lstlisting}[style=halfs,language=mlir]
(@\hl<2>{\lkw{scf.for} \%i = \%c0 \lkw{to} \%numRows \lkw{step} \%c1 \{}@)





  (@\hl<3>{\lcm{// steps 3 and 4 go here}}@)
}
\end{lstlisting}
    \end{column}
  \end{columns}
  \vspace{0.6em}
  \begin{itemize}
    \item<2-> Each \cd{ForOp::create} is one \cd{scf.for}.
    \item<3-> Ops created after \cd{setInsertionPointToStart} land
      \alert{inside} that loop's body.
    \item<4-> \cd{InsertionGuard} restores the insertion point when the scope
      ends. Leave it out and your ops land in the wrong block.
  \end{itemize}
\end{frame}
\note{Walk the three lines that repeat for every loop: create the op, move the
insertion point into its body, grab the induction variable.}

\begin{frame}[fragile]{Reference: \cd{MatVecToSCFPattern}, Step 3 of 5 --- Reduction Loop}
  \begin{columns}[t]
    \begin{column}{0.59\textwidth}
      \begin{center}\small\textbf{the C++}\end{center}
\begin{lstlisting}[style=halfs,language=C++,morekeywords={auto,Value}]
(@\hl<2>{\lkw{Value} init = memref::LoadOp::create(rewriter, loc, out,}@)
                                    (@\hl<2>{ValueRange\{i\});}@)
(@\hl<3>{\lkw{auto} redLoop = scf::ForOp::create(rewriter, loc, c0, numCols, c1,}@)
                                  (@\hl<3>{ValueRange\{init\});}@)
{
  (@\hl<4>{OpBuilder::InsertionGuard innerGuard(rewriter);}@)
  (@\hl<4>{rewriter.setInsertionPointToStart(redLoop.getBody());}@)
  (@\hl<4>{\lkw{Value} j = redLoop.getInductionVar();}@)
  (@\hl<4>{\lkw{Value} acc = redLoop.getRegionIterArgs()[0];}@)

  (@\hl<5>{\lkw{Value} matElem = memref::LoadOp::create(rewriter, loc, mat,}@)
                                         (@\hl<5>{ValueRange\{i, j\});}@)
  (@\hl<5>{\lkw{Value} vecElem = memref::LoadOp::create(rewriter, loc, vec,}@)
                                         (@\hl<5>{ValueRange\{j\});}@)
  (@\hl<6>{\lkw{Value} prod =}@)
      (@\hl<6>{arith::MulFOp::create(rewriter, loc, matElem, vecElem);}@)
  (@\hl<6>{\lkw{Value} sum = arith::AddFOp::create(rewriter, loc, acc, prod);}@)
  (@\hl<7>{scf::YieldOp::create(rewriter, loc, ValueRange\{sum\});}@)
}
\end{lstlisting}
    \end{column}
    \begin{column}{0.40\textwidth}
      \begin{center}\small\textbf{the IR it generates}\end{center}
\begin{lstlisting}[style=halfs,language=mlir]
(@\visible<2->{\hl<2>{\%init = \lkw{memref.load} \%out[\%i]}}@)

(@\visible<3->{\hl<3>{\%acc = \lkw{scf.for} }\hl<3-4>{\%j}\hl<3>{ = \%c0 \lkw{to} \%numCols \lkw{step} \%c1}}@)
    (@\visible<3->{\hl<3>{\lkw{iter\_args}(}\hl<3-4>{\%a}\hl<3>{ = \%init) -> (\lkw{f32}) \{}}@)






  (@\visible<5->{\hl<5>{\%m = \lkw{memref.load} \%mat[\%i, \%j]}}@)

  (@\visible<5->{\hl<5>{\%v = \lkw{memref.load} \%vec[\%j]}}@)

  (@\visible<6->{\hl<6>{\%prod = \lkw{arith.mulf} \%m, \%v : \lkw{f32}}}@)

  (@\visible<6->{\hl<6>{\%a2 = \lkw{arith.addf} \%a, \%prod : \lkw{f32}}}@)
  (@\visible<7->{\hl<7>{\lkw{scf.yield} \%a2 : \lkw{f32}}}@)
(@\visible<3->{\}}@)
\end{lstlisting}
    \end{column}
  \end{columns}
\end{frame}
\note{Six steps, one statement at a time:
1. The load of out[i]: the running total starts from what is already there,
   which is the +=.
2. ValueRange\{init\} makes the loop carry a value, starting at \%init.
3. Move into the body. j is \%j; acc is \%a, the carried value as seen from
   inside the body.
4. Read mat[i, j] and vec[j].
5. Multiply, and add to the carried value.
6. Yield the new value. The loop carries one, so the yield is explicit -- a
   loop with no iter\_args gets its scf.yield created for it.}

\begin{frame}[fragile]{Reference: \cd{MatVecToSCFPattern}, Step 4 of 5 --- Store}
  \begin{columns}[t]
    \begin{column}{0.62\textwidth}
      \begin{center}\small\textbf{the C++}\end{center}
\begin{lstlisting}[style=halfs,language=C++,morekeywords={LogicalResult,auto,Value,override}]
LogicalResult matchAndRewrite(MatVecOp op, ...) const override {
  // step 1: mat, numRows, numCols, c0, c1
  auto rowLoop = scf::ForOp::create(rewriter, loc, c0, numRows, c1);
  {
    // step 2: the guard, the insertion point, i
    Value init = memref::LoadOp::create(..., out, ValueRange{i});
    auto redLoop = scf::ForOp::create(..., ValueRange{init});
    {
      // step 3: the reduction body
    }
    (@\df{memref::StoreOp::create(rewriter, loc, redLoop.getResult(0), out,}@)
                            (@\df{ValueRange\{i\});}@)
  }
  rewriter.eraseOp(op);
  return success();
}
\end{lstlisting}
    \end{column}
    \begin{column}{0.37\textwidth}
      \begin{center}\small\textbf{the IR it generates}\end{center}
\begin{lstlisting}[style=halfs,language=mlir]
(@\strut@)

scf.for %i = %c0 to %numRows step %c1 {


  %init = memref.load %out[%i]
  %acc = scf.for %j ... {

    ...
  }
  (@\df{memref.store \%acc, \%out[\%i]}@)

}
\end{lstlisting}
    \end{column}
  \end{columns}
  \vspace{0.6em}
  \begin{itemize}
    \item The store takes the result of the \cd{scf.for},
      \cd{getResult(0)}: the value of the reduction, \cd{\%acc}.
  \end{itemize}
\end{frame}
\note{It goes after the reduction loop and inside the row loop. Store the
iter arg (getRegionIterArgs()[0]) instead and it builds, runs, and silently
gives the previous value of out.}

% =======================================================================
\begin{frame}[fragile]{Reference: \cd{MatVecToSCFPattern}, Step 5 of 5 --- Registration}
  \cd{LowerBlasToSCF.cpp}, the pass the pattern belongs to:
\begin{lstlisting}[language=C++,basicstyle=\ttfamily\scriptsize,morekeywords={void,override,struct,public}]
struct LowerBlasToSCFPass
    : (@\hl<2>{\lkw{public} blas::impl::LowerBlasToSCFPassBase<LowerBlasToSCFPass>}@) {
  void runOnOperation() override {
    RewritePatternSet patterns(&getContext());
    (@\hl<3>{patterns.add<MatVecToSCFPattern>(\&getContext());}@)
    // TODO(exercise): register your pattern too:
    //   patterns.add<MatMulToSCFPattern>(&getContext());

    (@\hl<4>{\lkw{if} (failed(applyPatternsGreedily(getOperation(), std::move(patterns))))} \{@)
      signalPassFailure();
      return;
    }
    ...
  }
};
\end{lstlisting}
  \vspace{0.3em}
  \begin{itemize}
    \item<2-> The pass's name, \cd{convert-blas-to-scf}, is declared in
      \cd{Passes.td}, in TableGen like the op, which generates this base
      class; the plugin registers the pass with \cd{mlir-opt}.
    \item<3-> \cd{patterns.add<...>} registers a pattern with the pass. Yours
      goes on the same line.
    \item<4-> \cd{applyPatternsGreedily} applies them to the IR the pass runs
      on --- next slide.
  \end{itemize}
\end{frame}

\begin{frame}{How Patterns Are Applied}
  \begin{itemize}
    \item<1-> \textbf{The driver applies the patterns.}
      \cd{applyPatternsGreedily} keeps applying the registered patterns until
      none of them matches any more.
    \item<2-> \textbf{Every IR change goes through the
      \cd{PatternRewriter}.} Creating, replacing and erasing ops all go through
      it, so the driver sees each change and can revisit what it affected. A
      pattern never edits the IR behind its back.
    \item<3-> \textbf{\cd{matchAndRewrite} reports what it did.} It returns
      \cd{success()} when it rewrote the op and \cd{failure()} when the op does
      not match --- and on \cd{failure()} it must leave the IR untouched.
  \end{itemize}
\end{frame}
\note{The driver also folds ops as it goes, which is why the pass's printed
output has constants where the pattern built memref.dim.}

% =======================================================================
\begin{frame}{Exercise 4 --- Write the Lowering}
  \begin{center}\large
    \cd{MatMulToSCFPattern} in\\\cd{blas-dialect/lib/LowerBlasToSCF.cpp}
  \end{center}
  \vspace{0.6em}
  \begin{itemize}
    \item<1-> One pass, \cd{convert-blas-to-scf}, one \alert{pattern per op}.
    \item<2-> The scaffolding is already there, commented out: operands,
      index constants, problem sizes, and the \cd{eraseOp} that retires the op.
    \item<3-> You write \alert{the loop nest} --- the same one as step 1, now
      built programmatically --- and then \alert{register the pattern}.
    \item<4-> Reference: \cd{MatVecToSCFPattern}, same file, complete.
  \end{itemize}
  \vspace{0.6em}
  \visible<5->{\begin{center}\color{scaccent}\large
    10 minutes to write it --- then we walk through the solution.
  \end{center}}
\end{frame}
\note{If you are behind schedule, this is the exercise to convert into a live
walkthrough: you type it, narrating, and they watch. The payoff here is
understanding the shape rather than producing it.

Remind them of the two rules from the reference: loop bodies are built by
moving the rewriter's insertion point into the loop's block, and an
InsertionGuard puts it back when the scope ends.}

% =======================================================================
\begin{frame}[fragile]{Solution --- \cd{MatMulToSCFPattern}, and What Changed from Matvec}
\begin{lstlisting}[language=C++,morekeywords={auto,Value},basicstyle=\ttfamily\scriptsize]
auto rowLoop = scf::ForOp::create(rewriter, loc, c0, M, c1);
{
  OpBuilder::InsertionGuard rowGuard(rewriter);
  rewriter.setInsertionPointToStart(rowLoop.getBody());
  Value i = rowLoop.getInductionVar();

  (@\df{auto colLoop = scf::ForOp::create(rewriter, loc, c0, N, c1);}@)  (@\only<2>{\color{scaccent}$\leftarrow$ one more loop: columns}@)
  (@\df{\{}@)
    (@\df{OpBuilder::InsertionGuard colGuard(rewriter);}@)
    (@\df{rewriter.setInsertionPointToStart(colLoop.getBody());}@)
    (@\df{Value j = colLoop.getInductionVar();}@)

    Value init = memref::LoadOp::create(rewriter, loc, out, ValueRange{i(@\df{, j}@)});
    auto redLoop = scf::ForOp::create(rewriter, loc, c0, K, c1, ValueRange{init});
    {
      // the guard, the insertion point, k and acc: as in matvec
      Value aElem = memref::LoadOp::create(rewriter, loc, a, ValueRange{i, k});
      Value bElem = memref::LoadOp::create(rewriter, loc, b, ValueRange{k(@\df{, j}@)});
      // prod, sum and the yield: as in matvec
    }
    memref::StoreOp::create(rewriter, loc, redLoop.getResult(0), out,
                            ValueRange{i(@\df{, j}@)});              (@\only<3>{\color{scaccent}$\leftarrow$ out and b take two indices}@)
  (@\df{\}}@)
}
\end{lstlisting}
\end{frame}
\note{This goes where the scaffolding says "your loop nest here". Everything
in black is the matvec pattern with the names changed (M, K, N for the sizes,
k for the reduction index). The orange is what is new: the column loop
wrapped around the reduction, and the second index on out and b.

It generates exactly the loop nest they wrote by hand in step 1.}

\begin{frame}[fragile]{Solution --- Register the Pattern}
\begin{lstlisting}[language=C++,basicstyle=\ttfamily\scriptsize,morekeywords={void,override}]
void runOnOperation() override {
  RewritePatternSet patterns(&getContext());
  patterns.add<MatVecToSCFPattern(@\df{, MatMulToSCFPattern}@)>(&getContext());
  ...
}
\end{lstlisting}
  \vspace{0.25em}
  \begin{itemize}
    \item<2-> A pattern that is never registered is never applied, and
      \cd{blas.matmul} survives the pass untouched.
    \item<3-> Which is the \alert{same symptom} as a pattern that returns
      \cd{failure()}. The pass catches both for you:
  \end{itemize}
  \vspace{0.3em}
  \visible<3->{\errbox{error: convert-blas-to-scf did not convert 'blas.matmul' -{}- is its pattern\\
    written and registered in blas-dialect/lib/LowerBlasToSCF.cpp?}}
\end{frame}

% =======================================================================
\begin{frame}{Recap}
  \begin{center}
  \begin{tikzpicture}[
      box/.style={draw=scblue,fill=scwash,rounded corners=2pt,inner sep=5pt,
                  font=\ttfamily\scriptsize,minimum height=2.4em,align=center},
      lbl/.style={font=\scriptsize\itshape,color=scgrey},
      arr/.style={-{Stealth[length=2mm]},draw=scgrey,thick}]

    \node[box] (scf) {scf.for nest};
    \node[box,right=1.6cm of scf] (llvm) {llvm dialect};
    \node[box,right=1.9cm of llvm] (run) {JIT + run};
    \draw[arr] (scf) -- node[lbl,above=2pt]{7 passes} (llvm);
    \draw[arr] (llvm) -- node[lbl,below=3pt,font=\tiny\itshape]{ExecutionEngine} (run);

    \node[box,below=1.2cm of scf,draw=scaccent] (op) {blas.matmul};
    \draw[arr,draw=scaccent] (op) -- node[lbl,right,align=left]
      {convert-blas-to-scf\\\emph{you wrote this}} (scf);

    \node[lbl,below=0.35cm of run,align=center,text width=3.4cm]
      {the same seven passes\\as step 1, same answer};
  \end{tikzpicture}
  \end{center}
  \vspace{0.2em}
  \begin{itemize}
    \item<2-> \cd{convert-blas-to-scf} turns \cd{blas.matmul} into the
      \cd{scf.for} nest you wrote by hand in step 1.
    \item<3-> An \alert{op} --- operands, types, assembly syntax --- from
      eight lines of ODS.
    \item<4-> A \alert{verifier}, so bad IR is rejected with your own error
      message.
    \item<5-> A \alert{pass} that rewrites your op into standard dialects.
  \end{itemize}
  \vspace{0.3em}
  \visible<6->{\begin{center}\large
    That is a complete dialect: \alert{definition}, \alert{validation},
    \alert{meaning}.\end{center}}
  \visible<7->{\begin{center}\small\color{scgrey}
    Everything upstream MLIR does --- \cd{linalg}, \cd{vector}, \cd{gpu} ---
    is this, at scale.\end{center}}
\end{frame}
\note{From convert-blas-to-scf onwards the rest of the compilation is what they
did in step 1: the same seven passes, the same JIT, the same check against
numpy. Their dialect only had to lower to something the rest of the stack
already understands.

End here. Part 2 has its own deck, notebook and runtime.}

\end{document}
